\documentclass[11pt]{article}
\usepackage[margin=1in]{geometry}
\usepackage[T1]{fontenc}
\usepackage[utf8]{inputenc}
\usepackage{lmodern}
\usepackage{booktabs}
\usepackage{tabularx}
\usepackage{array}
\usepackage{amsmath}
\usepackage{microtype}
\usepackage{enumitem}
\usepackage[hidelinks]{hyperref}
\usepackage[numbers,sort&compress]{natbib}

\newcolumntype{L}{>{\raggedright\arraybackslash}X}
\setlist{itemsep=2pt,topsep=4pt}

\title{Grounding Documents for Long-Running Tabletop Campaigns:\\
A Practitioner's Survey of Twelve Approaches,\\
from LLM Extraction to Authored Structure}

\author{Kostadis Roussos\\
Independent researcher\\
\texttt{kostadis@gmail.com}}

\date{October 2026}

\begin{document}
\maketitle

\begin{abstract}
A game master who uses language models to prepare sessions of a long-running tabletop role-playing
campaign needs \emph{grounding documents}: a current world state, campaign state, party record,
planning notes and per-character dossiers, distilled from dozens of recorded sessions and fed to the
models as context. Producing these documents is easy to start and hard to get right, because the
details that matter are exactly the ones generative models get wrong: who did what, in what order,
which name denotes which entity, and whether something happened at all. Over seven months
(March--October 2026) and six campaigns of up to 70 chapters, one practitioner built, measured and
retired twelve distinct approaches to the problem. We reconstruct them from the git history,
design documents and experiment write-ups of the repositories involved, and organise them by two
axes: the \emph{substrate} a pipeline reads (narrated chapter prose, atomic extracted facts,
retrieval indexes, or the structured session summaries the GM already reviews) and \emph{who makes
the precision decisions} (ordering, attribution, identity, scope). Approaches that let a model make
those decisions failed in ways that were silent and that compounded: a consolidated extraction shrank
a planning document from 513 to 109 lines; an atomic-fact ensemble reported a slain antagonist as
alive, and 98\% of its 15{,}465 facts were seen by only one sample, so agreement could not serve
as a filter; a knowledge-graph library stored 1.18 characters of generated text per character of
source and answered ``what was the most recent session'' with session~8 of~11. The approach that the
GM finally promoted parses the scene order, entity sections and memorable moments already authored in
reviewed session summaries, with no model call (67 summaries, 409 scenes, 868 dossiers,
byte-identical across runs), and uses a model only to render the final prose; it also flagged 51
claims in the GM's own tracking files as unsupported by the record. We distil the history into ten
cross-cutting findings and a set of design rules for grounding long-horizon generative work.
\end{abstract}

\section{Introduction}

Tabletop role-playing campaigns are long. A weekly game runs for years, and the record of what the
players have done---who they met, what they promised, which villain is dead and which only appears to
be---is the substrate from which the game master (GM) prepares every subsequent session. When the GM
asks a language model to help with preparation (draft an encounter, voice a returning character,
check a plot thread), the model needs that record as context. Feeding it hundreds of thousands of
words of transcript is impractical and, past a point, counter-productive~\cite{liu2024lost}. The
practical answer is a small set of \emph{grounding documents}: compact, current, authoritative
summaries of the world, the campaign, the party and the plans, together with per-entity dossiers.

This paper surveys how one practitioner tried to build those documents. The work was done in an
open-source toolkit for campaign management (CampaignGenerator~\cite{campaigngen}), a private
repository of campaign data for six campaigns, a public repository of local-model experiments on two
NVIDIA DGX Spark machines~\cite{dgxfun}, and a number of local-only scratch directories. Several
approaches are still live, several survive only in history, and some never reached version control.
The aim is to make that history usable: to describe each approach precisely enough that it can be
compared with the others, to record the measured reason each one was kept or dropped, and to extract
what generalises.

The problem looks like summarisation, and the first approaches treated it that way. It is better
understood as \emph{record keeping under precision constraints}. A grounding document is consumed by
another model that cannot tell a confident error from a fact. Errors in four kinds of decision are
both likely and expensive: \emph{ordering} (what happened before what), \emph{attribution} (who did
it, to whom), \emph{identity} (which names denote the same entity) and \emph{scope} (what belongs in
this document, and what actually happened as opposed to what was planned). The toolkit's governing
rule, stated before any of the work surveyed here, is that a model may draft but these
\emph{precision decisions} require a human checkpoint before they feed another model. Much of the
history below is the gradual discovery of where that rule had been violated without anyone noticing,
and of how to remove the model from those decisions rather than merely review its output.

\paragraph{Contributions.}
\begin{itemize}
  \item A catalogue of twelve approaches to producing grounding documents for long-running campaigns,
        reconstructed from primary sources, with status, lineage and measured outcomes
        (Sections~\ref{sec:taxonomy}--\ref{sec:approaches}).
  \item A two-axis taxonomy (substrate $\times$ locus of precision decisions) under which the
        successes and failures separate cleanly (Section~\ref{sec:taxonomy}).
  \item Ten cross-cutting findings, each with a measurement, that apply beyond tabletop games to any
        long-horizon generative pipeline that must keep a record (Section~\ref{sec:findings}).
\end{itemize}

\section{Background}
\label{sec:background}

\subsection{The campaigns and their record}
Six campaigns contributed evidence. They are adaptations of published adventures: \emph{Out of the
Abyss} (70 chapters at the time of writing), a hybrid of \emph{Dragon of Icespire Peak} and
\emph{Lost Mine of Phandelver} (53 sessions), \emph{Storm King's Thunder}, \emph{Temple of Elemental
Evil}, an organised-play adventure set in Hillsfar, and a second \emph{Lost Mine} table (12
sessions). Each session is recorded and transcribed. A post-session pipeline produces, in order: a
reviewed transcript (spelling and speaker attribution corrected by the GM); a structured
\emph{session summary} with fixed sections (\texttt{Summary}, \texttt{Memorable Moments}, an ordered
\texttt{Scenes} list, and entity sections for NPCs, locations, items and spells); scene extractions;
and finally a narrated prose chapter. The narrated chapters are concatenated into a campaign
\emph{bible} of several megabytes.

\subsection{The grounding documents}
Four documents form the core: \emph{world state} (the setting as it now stands), \emph{campaign
state} (where the story is, open threads, recent events), \emph{party} (characters, resources, arcs)
and \emph{planning} (NPC dossiers, factions, threat tracking). They are read by the toolkit's
preparation agents and by interactive assistants. A per-campaign \emph{entity registry} records
canonical names, types and approved aliases for about a thousand entities per campaign; it is the
toolkit's authority on identity and spelling and, as Section~\ref{sec:findings} argues, it is
deliberately a names layer, not a state layer.

\subsection{Related work}
Retrieval-augmented generation~\cite{lewis2020rag} and graph-structured retrieval~\cite{edge2024graphrag}
are the standard alternatives to static context documents; tiered memory systems for agents
\cite{packer2023memgpt,park2023generative} maintain summaries and reflections that are themselves
model-written. Book-length summarisation~\cite{chang2024booookscore} measures the coherence errors
that accumulate when a model compresses a long narrative incrementally. Work on language models in
\emph{Dungeons \& Dragons} has used actual-play logs and structured game state as training and
evaluation data~\cite{callisonburch2022dnd,zhu2023fireball} and has built assistants for the
dungeon master at the table~\cite{zhu2023calypso}. None of these addresses the maintenance problem
surveyed here: keeping a reviewed, provenance-bearing record of a campaign current over months, under
a rule that a human, not a model, settles ordering, attribution and identity. The hallucination
literature~\cite{ji2023survey} catalogues the failure modes; this survey records which pipeline shapes
expose a practitioner to them and which do not.

\section{Survey method}
\label{sec:method}
The sources are (i) the full git history of the toolkit repository, including all 28 local and remote
branches, deleted files and merged pull requests; (ii) the campaign-data repository (two clones,
about 130 branches); (iii) the experiment repository and eleven further repositories or forks, of
which only the practitioner's own commits were considered; (iv) local directories outside version
control; and (v) the design documents and experiment write-ups committed alongside the code, which
record measurements and kill decisions at the time they were made. Three independent passes
catalogued the repositories, and their reports were reconciled against the primary documents; every
number in this paper is quoted from a design document, experiment write-up or commit message.

An approach is counted as \emph{distinct} if it changes the substrate, the locus of a precision
decision, or the output form; a parameter change within an approach is not counted. Status is one
of: \emph{live} (on the main branch and used for at least one campaign's current documents);
\emph{dormant} (on main, not used); \emph{history} (deleted, recoverable from git);
\emph{proposal} (designed and measured, not built); \emph{local} (exists only outside version
control).

\section{A taxonomy}
\label{sec:taxonomy}

Two axes separate the approaches.

\paragraph{Substrate.} What the pipeline reads: (P) narrated \emph{prose}, the bible chapters;
(F) atomic \emph{facts} extracted from prose by a model; (R) a \emph{retrieval} index over some
corpus, queried at preparation time; (S) the \emph{structured} session summaries that the GM already
reviews as part of the post-session pipeline. The narrated prose is downstream of the summaries: a
pipeline that reads prose to recover structure is undoing a rendering step.

\paragraph{Locus of precision decisions.} Who decides ordering, attribution, identity and scope:
a model (M); deterministic code over declared structure (D); or the GM through a recorded ruling (G).
Most approaches mix loci; Table~\ref{tab:catalogue} gives the dominant one for the step that
determines what the document asserts.

\begin{table}[t]
\centering
\small
\caption{The twelve approaches. Substrate: P prose, F atomic facts, R retrieval index,
S structured summaries. Locus: who makes the ordering/attribution/identity/scope decisions---M model,
D deterministic code, G GM ruling.}
\label{tab:catalogue}
\begin{tabularx}{\textwidth}{@{}r L c c l L@{}}
\toprule
\# & Approach & Subst. & Locus & Active & Status \\
\midrule
1  & Per-document extract-and-synthesise scripts         & P & M     & Mar--Aug & live (3 campaigns), dormant elsewhere \\
2  & Shared chapter-extract consolidation                & P & M     & Apr      & history (killed) \\
3  & Hand curation and dossier merge                     & P & G     & Apr--    & live as stopgap \\
4  & Local atomic-fact ensemble                          & F & M     & May--Sep & live, no longer feeds OOTA \\
5  & Ensemble repairs: recency, offsets, coverage        & F & D/G   & Jul      & live (within 4) \\
6  & Two-stage finder/claim extraction                   & F & D     & Sep      & proposal \\
7  & Entity-level typing and ruling replay               & F & D/G   & Oct      & proposal, replayed \\
8  & Model work over ensemble dossiers                   & F & M     & Jul--Aug & local \\
9  & State stores and per-section projections            & F & D/G   & Jul--Aug & dormant \\
10 & Knowledge-graph and semantic retrieval              & R & M     & Jul--Sep & dormant (graph); live (search) \\
11 & Event-ordinal graph over authored scenes            & S & D/G   & Sep      & proposal, local outputs \\
12 & Summary-native grounding                            & S & D/G   & Oct      & live (OOTA) \\
\bottomrule
\end{tabularx}
\end{table}

The table already shows the result of the survey: the approaches that a GM promoted, or that measured
well, are those in which the precision decisions moved out of the model, and the two that read
structured summaries (11 and 12) are the only ones in which the model is never asked to recover
an order or an identity.

\section{The approaches}
\label{sec:approaches}

\subsection{Family P/M: extract from prose, synthesise with a model}

\paragraph{1. Per-document scripts (March 2026).}
Four scripts---one per grounding document---each split the bible into chapter-aligned chunks of up
to about 60{,}000 characters, ran a model extraction per chunk into their own extraction directory,
and ran a second model call to synthesise the document from those extractions. A
\texttt{-{}-build-dossiers} mode of the planning script wrote per-NPC dossiers. The human role was
review after the fact. The approach is simple and the documents read well, and three campaigns still
use it. Its failures are of staleness and invention rather than style. An issue note from July
records that ``the extraction layer has Chapter~62. The synthesis layer did not consume it'': all four
documents were stamped Chapter~62 and described Chapter~61. In another campaign the synthesis
introduced a non-existent NPC that had to be removed by hand. Each script also re-reads the same
input, which motivated approach~2.

\paragraph{2. Shared chapter extract (killed, April 2026).}
The consolidation replaced three narrow extractions with one rich extraction per chapter, in an
eight-section schema, consumed by three synthesisers. At the extraction layer it looked better: the
schema held across chapters of 2.7K--54K characters and one file per chapter was easier to review. At
the synthesis layer every document regressed. World state fell from 268 to 193 lines and lost its
party and items sections; campaign state's NPC table fell from 40 to 23 rows; planning fell from 513
to 109 lines, with 17 NPC dossiers collapsing to 5 and threat-tracker values ``\emph{guessed}
(`6--8 estimated')'' instead of read. Outputs were well under the token cap; this was depth loss, not
truncation. The kill note's diagnosis is general: ``Consolidating $N$ narrow LLM extracts into one
broad extract is not free. The narrowness \emph{is} the depth.'' It also named the principle that
recurs throughout this survey: ``the rough extraction pass is the ceiling, not the floor.''

\paragraph{3. Hand curation.}
Underneath every generated approach sits human repair: merging duplicate dossiers (one pass merged 54
files into 25), sidecar files of new notes, candidate files for the party document so that a
regeneration would not overwrite curated text, and hand fold-ins that carry a document forward from
the latest session summaries without regenerating it. A commit message records the cost: dossiers
``are generated and \emph{then hand-curated}\ldots\ Regenerating produces the machine output, not the
curated state.'' Hand fold-in is the current state of two campaigns' documents, and it coexists
uneasily with a stated rule that generated documents are not hand-edited.

\subsection{Family F: atomic facts}

\paragraph{4. Local atomic-fact ensemble (May--September 2026).}
To extract once and share the result, extraction moved to local hardware: five extraction
\emph{lenses} (prompts emphasising different material) times three samples, served from an
80B-parameter open model across two DGX Spark machines, emit atomic facts with a type and a subject.
A deterministic merge groups facts by (type, normalised subject), or optionally by embedding
similarity. A bundling step writes per-entity dossiers, with a listing mode documented as ``the human
checkpoint on what will be aggregated''. GM skills review alias and type merges. One model call per
document then synthesises the grounding documents from selected dossiers. The first full run
(1~June) produced 22{,}240 facts over 56 chapters and 703 dossiers; a parallel set of documents was
written but, per its own commit, they were ``not drop-in replacements\ldots\ campaign\_state is a
derivation-of-derivation,'' and that set was later deleted. One smaller campaign promoted ensemble drafts
after GM review; the largest campaign never did.

A July investigation of the stale-chapter report found that the approach's deepest problem was not a
bug. ``Atomization strips the syntax that carries attribution.'' A sentence that identifies a dead
antagonist is unambiguous in its paragraph and unrecoverable as a fragment; in the case that drove the
investigation, three facts were extracted from one quote, two correct and one inverting the roles,
and because bundling is by subject the inverted fact became the \emph{only} fact in the antagonist's
dossier. Synthesis then amplified it to ``Current status: Alive''. The obvious filter does not exist:
98\% of facts (15{,}220 of 15{,}465) were produced by a single sample, because the lenses extract
different material by design. Embedding similarity does not help either, since two statements that
disagree about who died are far apart by construction (cosine 0.85 between the correct and inverted
fact, below the merge threshold), and a contradiction detector found 15 candidate groups, 14 of
them benign.

\paragraph{5. Repairs within the ensemble (July 2026).}
The investigation's through-line was: ``The pipeline computes a signal, serialises it, and then
ignores it---forcing a later stage to guess what it was already told.'' Four instances were fixed in a
day. A dossier frequency floor ranked by fact count, which ``deleted the newest chapter'' because
new entities have the fewest facts; a recency window restored all seven entities of the latest chapter.
The embedding merge was unreachable from the interface, so ``62 chapters silently merged with the
weaker algorithm''. Quote positions were reduced to a boolean, which lost within-chapter order;
stamping offsets made the inverted fact visibly 800 characters after the scene that refutes it.
Entities attested only through other entities' facts were invisible; a coverage report surfaced 30,
including one with zero facts of its own and 34 mentions that had no dossier at all. A chapter--summary
mapping and a per-chapter narrative pass were added, both gated on explicit GM approval. These repairs
made errors \emph{legible}; they did not make the facts correct.

\paragraph{6. Two-stage extraction (proposal, September 2026).}
Replaying the shipped merge over a fresh corpus of 13{,}731 facts measured what the prose comparator
costs: 2{,}384 pairs of rows whose source spans overlap but were kept separate (821 with
byte-identical quotes), and 43 merged rows whose members sit on disjoint spans---in one, three of a
chapter's own date headings fused into a row asserting one date. The proposal splits extraction into a
finder stage that produces spans and a claim stage that cites them, so the span becomes the merge key.
An open-source grounded-extraction library was evaluated as the finder and rejected: on facts it was
not stable across chapter size, falling to 6.2\% source coverage on a 14K-character chapter against
75.6\% for the ensemble union, as extractions degenerated into bare names.

\paragraph{7. Entity-level typing and ruling replay (October 2026).}
Typing each fact independently splits an entity across types (an NPC in one chapter, a monster in the
next); the GM had resolved 109 such splits by hand in one campaign. A companion study~\cite{roussos2026decision}
found that re-typing facts with a non-generative decision model ``creates roughly as many splits as it
repairs.'' The replacement types once per entity through a deterministic chain: GM filing conventions,
then prior GM merge rulings, then the registry type, then a 60\% majority vote, then a GM queue. The
model-free vote matched the GM's filing 32 times in 33 where it decided. The proposal's sharper
conclusion was that the pipeline ``asks the LLM for precision decisions and never reads the GM's
rulings back'': rulings were recorded and then ignored on the next run.

\paragraph{8. Model work layered on ensemble dossiers (July--August 2026).}
Several local jobs, run by a locally served model through an agent harness and kept outside version
control, worked on the ensemble's dossiers: appending matching facts to each dossier; answering each
dossier's uncertainty bullets from the cited chapters (327 dossiers: 158 resolved, 140 partial, 11
unresolved, 18 malformed outputs; four dossiers found contradicted by the chapters); per-entity chapter
scans and an inferred calendar timeline; and a cross-campaign ``scholar persona'' wiki with a sampling
fact-checker. None was promoted. Their main yield was evidence that unreviewed model digests drift on
names and that dossiers built from atomic facts carry contradictions that a later pass can find but
not settle.

\paragraph{9. State stores and projections (July--August 2026).}
This approach changed the output form rather than the extraction. Durable stores---a deduplicated
per-chapter event spine, a registry of GM-ratified plot threads with model-proposed candidates, and a
GM-approved chapter--session map---feed per-section renders, and a document is the assembly of its
sections. A build without a model backend makes ``zero model calls; a re-run with unchanged stores
rebuilds nothing.'' Planning is layered by certainty, from ratified threads down to explicitly labelled
non-canon speculation. The design is sound; its inputs were not. The event spine was built from
ensemble facts, and near-duplicate phrasings of one event survived its 60-character deduplication key.
In the one campaign that used it, four threads had been ratified and 148 proposals were pending when work
stopped, and that campaign's live world state returned to hand fold-in.

\subsection{Family R: retrieval instead of documents}

\paragraph{10. Knowledge graphs and semantic search (July--September 2026).}
A knowledge-graph library~\cite{cognee} was used to ingest grounding documents and dossiers for one
campaign and session summaries for another, with a locally served model building the graph. The first
attempt never completed in about six runs; a later retest completed in 27 minutes and answered five
evaluation queries correctly, but the graph was never adopted, because the structure was written by a
model with no review. The September investigation explains why that matters. The library stored
2{,}468{,}434 characters of generated text for 2{,}093{,}539 characters of source (1.18$\times$),
including paraphrase drift, facts imported from the published adventure that never happened at the
table (327 edges about a castle the party has never found), and a false identity merge of a villain
with his master. A stripped prompt did not stop it: with structured output, ``the schema is the
stronger instruction.'' Asked for the most recent session of an eleven-session campaign, the graph
returned Chapter~8, a vector search tool returned session~6 and a semantic memory store~\cite{mempalace}
returned session~8. ``Every semantic tool failed the time question.'' The verbatim tools never
fabricated; their failures were in ranking. A semantic search skill over the summaries remains live as
a \emph{lookup} aid alongside the documents, with filename ordering and an explicit absence check added
because ``vector search structurally cannot'' supply them.

\subsection{Family S: authored structure}

\paragraph{11. Event-ordinal graph (September 2026).}
The September investigation's conclusion was that ``the sequence of events is already authored'':
filenames give session order, each summary's \texttt{Scenes} section gives scene order, and its bullets
give beat order. ``Time is the event ordinal, not the calendar.'' A graph keyed on that ordinal, with
state snapshots from the summaries' entity sections and identity taken from the GM-reviewed alias file
(``the one precision decision''), was built in about a second for twelve sessions (72 events, 936 beats,
253 snapshots), answered every time question the semantic tools failed, and was reproduced
byte-identically by three independent extractions. Retargeted to a 53-session campaign it ran in
0.8~s (356 events, 5{,}255 beats). Identity resolution through the alias file raised one NPC's
mentions from 4 to 15 and, unlike the model-built graph, did not merge the villain with his master.
The investigation also observed that the alias file mixed surface-form equivalences with secret
identities known only to the GM, and that an identity revelation is itself an event. Deterministic
state documents were generated from the graph by exact search and copy, with no model. They were
never promoted, and the proposal to replace the per-document campaign-state script was not built.

\paragraph{12. Summary-native grounding (October 2026).}
The approach now in use for the largest campaign reached the same design independently, starting from a
failed attempt to run the ensemble over summaries instead of prose (a four-lens extraction projected at
7--10 hours on the available hardware). It reads a directory of structured summaries and nothing else.
Validation runs first, collects every problem across the corpus, and refuses to proceed on any error in
range: filename number, title number and scene identifiers must agree, and a file without a
\texttt{Scenes} section is never treated as one undifferentiated block. A deterministic parse then
produces a chronology of every scene, the memorable-moments sections preserved character for
character, and one dossier per entity with the source file and scene of every observation. Headings
are grouped only when identical within a category or listed as exact same-type aliases in the
registry. Likely duplicates (similarity $\ge 0.88$, or a name with and without a parenthetical
qualifier) are \emph{listed, never merged}: ``No merge decision MAY be made by a model or a similarity
score''---the GM corrects the summary, or records the pair as distinct in a rulings file the pipeline
only reads. Every deterministic artifact is byte-identical across runs given identical inputs.

Only then does a model write: world state, campaign state, party and planning drafts, from dossiers
selected by recency and recurrence and from the full chronology. Tracking and planning files are given
to campaign-state synthesis as \emph{questions to audit}, not as evidence, and any claim the summaries
do not support is labelled rather than narrated. A draft feeds the next synthesis only when the GM names
it as reviewed; drafts never overwrite live documents; exact prompts, input digests and selection
decisions are retained. On the full corpus the parse covered 67 of 67 summaries, 409 scenes and 868
dossiers; the drafts extended the live documents from Chapter~63 to Chapter~70; and the campaign-state
audit marked 51 tracking claims as not found in the summaries---planned or module-default events that
never happened at the table. The GM judged the results ``much better than the other ones we have
built'' and promoted all four.

\subsection{Cross-cutting: the entity registry}
Every approach from 4 onward depends on the entity registry for identity. It was built from module
inventories through a model pass, a deterministic merge and GM rulings, imported by tooling rather than
edited by hand (954 entities in one campaign, ``no manual YAML surgery''), and later cut from 990 to 922
entries in another by removing misspellings that had been recorded as aliases. The rule that emerged is
that the registry records names, types and aliases, never state or history.

\section{Findings}
\label{sec:findings}

\paragraph{F1. The extraction pass is the ceiling, not the floor.}
An impressive first-pass extraction predicts nothing about the synthesis built on it (approach~2), and
a synthesis cannot recover detail that the extraction never wrote down at the needed granularity.

\paragraph{F2. Atomisation destroys attribution.}
Decontextualised facts are what make deduplication and counting possible, and they are also why ``who
did what to whom'' cannot be recovered downstream by clustering, voting or embedding (approach~4).

\paragraph{F3. Agreement is not available as a filter.}
Ensembles are usually justified by agreement. With lenses designed to extract different material, 98\%
of facts had a single supporting sample, and embeddings measure paraphrase, not contradiction.

\paragraph{F4. Computed signals get discarded, and the system then guesses.}
Chapter ranges, merge method, quote positions and mention counts were all computed and then ignored by
the next stage (approach~5). The repair in each case was to carry the signal forward, not to add a model.

\paragraph{F5. Silent success is the dominant failure mode.}
A library ingested a file path as text and reported success; a merge fell back to a weaker algorithm with
a message that read as configuration; a reasoning model spent its whole completion budget thinking and
returned no content; three tools reported success over work they had not done. Instruments need positive
controls.

\paragraph{F6. Time is an ordinal, and semantic retrieval cannot find it.}
Every semantic tool failed ``most recent session''. Calendar dates are sparse and inconsistently
formatted; the authored scene order is always present (approaches~10, 11).

\paragraph{F7. Generative structure invents the record.}
Graph-building and summarising layers import module canon as history, drift in paraphrase and merge
identities, and store the result as fact (approach~10). Verbatim tools fail by omission, which a
reviewer can detect; generative layers fail by commission, which a reviewer must catch.

\paragraph{F8. Identity is the precision decision.}
Which names denote the same entity is the one decision that cannot be read off the text, so it belongs
to the GM, recorded once and read back every run (approaches~7, 11, 12). Surface-form equivalence and
secret identity are different decisions and need different scopes.

\paragraph{F9. Plans are not evidence.}
Tracking files, planning notes and module text describe what was expected. Treating them as audit
questions exposed 51 claims that never happened (approach~12).

\paragraph{F10. Do not undo a rendering step.}
The narrated bible is rendered from the structured summaries. Re-extracting structure from it is a
lossy round trip---``structure $\to$ render to prose $\to$ re-extract structure from the prose''---and
the dimension it loses (order, attribution) is the one the extractors then spend tokens guessing. The
cheapest, fastest and most accurate approach read the structure that already existed.

\section{Discussion}

\paragraph{Where the model belongs.}
Across the twelve approaches the model's role shrank from extractor, structurer and renderer
(approaches~1--2) to renderer only (approach~12), and each step of that shrinkage was forced by a
measured failure, not by principle. The final division is: authored structure supplies order; the
registry and recorded GM rulings supply identity; deterministic code supplies grouping, selection and
provenance; the model supplies prose. The human checkpoint did not move---it was always the GM---but
it moved \emph{earlier}, from reviewing generated documents to reviewing the summaries and rulings the
documents are computed from, where a correction is made once and persists.

\paragraph{Rulings must be inputs.}
Several approaches asked the GM for a decision, recorded it, and asked again on the next run. A
ruling store that every stage reads---conventions, merges, aliases, not-a-duplicate pairs---is the
mechanism that keeps a checkpoint from becoming a recurring chore that an exhausted reviewer
eventually rubber-stamps.

\paragraph{What remains open.}
Summary-native grounding moves the precision problem upstream; it does not eliminate it. The summaries
are themselves produced with model assistance and reviewed by the GM, and their quality bounds
everything downstream; chapters without recordings need summaries written from the narrated prose,
reintroducing a rendering step. The store-and-projection design (approach~9) is the natural way to make
summary-native outputs incremental and section-addressable, but it was built over ensemble facts and
has not been retargeted. Three campaigns still run approach~1, and migrating them requires their
summaries to conform to the structured schema. Finally, the history itself was at risk: the prototype
that led to approach~12 was built in a temporary worktree whose code was never committed, and the
event-ordinal outputs and several local experiments exist only on one machine. A practitioner running
many exploratory approaches should commit the evidence even when the code is disposable.

\paragraph{Cost and local hardware.}
The local ensemble was attractive because extraction on owned hardware is nearly free per run. It
was not free in wall-clock time (hours per corpus, 7--10 hours projected for one summary run) or in
review time (tens of thousands of facts no one reads). The deterministic approaches run in seconds.
The decisive saving was not cheaper extraction but not extracting.

\section{Limitations}
This is a retrospective survey of one practitioner's work, reconstructed by an AI agent from
repositories and documents written during the work; it inherits their framing. The approaches were not
evaluated on a common benchmark: measurements were taken on different campaigns, corpus versions and
models, and comparisons across approaches are qualitative. The GM's judgement of the final documents is
a single expert's assessment, not a blind comparison. Status labels for campaigns not examined in depth
(in particular which method produced one campaign's latest refresh) are uncertain. Local-only work was
assessed from its outputs and scripts, without commit history. Summary-native grounding has been
promoted for one campaign; its behaviour on campaigns with less regular summaries is untested.

\section{Conclusion}
Twelve attempts to build grounding documents for long-running campaigns, spread over seven months and
six campaigns, converge on a simple conclusion: the record a GM needs is already written, in the
structured session summaries the GM reviews after every session, and the job of tooling is to
read that structure faithfully, keep identity decisions with the GM, and let a model do only what it is
good at, which is writing. The approaches that asked a model to recover order, attribution or identity
from prose or fragments failed silently and compounded their errors; the approaches that read authored
order and recorded rulings were faster, reproducible and better, by the only judge whose opinion
counted.

\section*{Data and code availability}
The toolkit, including the design documents, experiment write-ups and the summary-native pipeline,
is public~\cite{campaigngen}, as are the local-model experiments~\cite{dgxfun}. The campaign data is
not released because it quotes play sessions verbatim and names real people.

\section*{Use of AI assistance}
The approaches surveyed were largely implemented by AI coding agents (Claude, Anthropic, and others)
under the author's direction. This survey was compiled by an AI agent from the repositories' history and
documents, and drafted by it; the author made the decisions the paper records and takes full
responsibility for the content.

\bibliographystyle{plainnat}
\begin{thebibliography}{99}

\bibitem[Roussos(2026a)]{campaigngen}
K.~Roussos.
\newblock CampaignGenerator: tooling for tabletop campaign management with language models.
\newblock \url{https://github.com/kostadis/CampaignGenerator}, 2026.

\bibitem[Roussos(2026b)]{dgxfun}
K.~Roussos.
\newblock dgx-fun: experiments with local LLM serving on NVIDIA DGX Spark.
\newblock \url{https://github.com/kostadis/dgx-fun}, 2026.

\bibitem[Roussos(2026c)]{roussos2026decision}
K.~Roussos.
\newblock Decision models at the game table: where non-generative classifiers help and fail in LLM
  tooling for tabletop role-playing campaigns.
\newblock Manuscript, October 2026.

\bibitem[Topoteretes(2026)]{cognee}
Topoteretes.
\newblock Cognee: memory for AI agents.
\newblock \url{https://github.com/topoteretes/cognee}, 2026.

\bibitem[MemPalace(2026)]{mempalace}
MemPalace contributors.
\newblock MemPalace: a semantic memory palace for agents.
\newblock Open-source software, 2026.

\bibitem[Lewis et~al.(2020)]{lewis2020rag}
P.~Lewis, E.~Perez, A.~Piktus, F.~Petroni, V.~Karpukhin, N.~Goyal, H.~K{\"u}ttler, M.~Lewis,
  W.~Yih, T.~Rockt{\"a}schel, S.~Riedel, and D.~Kiela.
\newblock Retrieval-augmented generation for knowledge-intensive {NLP} tasks.
\newblock In \emph{Advances in Neural Information Processing Systems}, 2020.

\bibitem[Edge et~al.(2024)]{edge2024graphrag}
D.~Edge, H.~Trinh, N.~Cheng, J.~Bradley, A.~Chao, A.~Mody, S.~Truitt, and J.~Larson.
\newblock From local to global: a graph {RAG} approach to query-focused summarization.
\newblock arXiv:2404.16130, 2024.

\bibitem[Packer et~al.(2023)]{packer2023memgpt}
C.~Packer, S.~Wooders, K.~Lin, V.~Fang, S.~G. Patil, I.~Stoica, and J.~E. Gonzalez.
\newblock {MemGPT}: towards {LLMs} as operating systems.
\newblock arXiv:2310.08560, 2023.

\bibitem[Park et~al.(2023)]{park2023generative}
J.~S. Park, J.~C. O'Brien, C.~J. Cai, M.~R. Morris, P.~Liang, and M.~S. Bernstein.
\newblock Generative agents: interactive simulacra of human behavior.
\newblock In \emph{Proceedings of UIST}, 2023.

\bibitem[Chang et~al.(2024)]{chang2024booookscore}
Y.~Chang, K.~Lo, T.~Goyal, and M.~Iyyer.
\newblock {BooookScore}: a systematic exploration of book-length summarization in the era of {LLMs}.
\newblock In \emph{International Conference on Learning Representations}, 2024.

\bibitem[Liu et~al.(2024)]{liu2024lost}
N.~F. Liu, K.~Lin, J.~Hewitt, A.~Paranjape, M.~Bevilacqua, F.~Petroni, and P.~Liang.
\newblock Lost in the middle: how language models use long contexts.
\newblock \emph{Transactions of the Association for Computational Linguistics}, 12:157--173, 2024.

\bibitem[Callison-Burch et~al.(2022)]{callisonburch2022dnd}
C.~Callison-Burch, G.~S. Tomar, L.~J. Martin, D.~Ippolito, S.~Bailis, and D.~Reitter.
\newblock Dungeons and dragons as a dialog challenge for artificial intelligence.
\newblock In \emph{Proceedings of EMNLP}, 2022.

\bibitem[Zhu et~al.(2023a)]{zhu2023fireball}
A.~Zhu, K.~Aggarwal, A.~Feng, L.~J. Martin, and C.~Callison-Burch.
\newblock {FIREBALL}: a dataset of dungeons and dragons actual-play with structured game state
  information.
\newblock In \emph{Proceedings of ACL}, 2023.

\bibitem[Zhu et~al.(2023b)]{zhu2023calypso}
A.~Zhu, L.~J. Martin, A.~Head, and C.~Callison-Burch.
\newblock {CALYPSO}: {LLMs} as dungeon masters' assistants.
\newblock In \emph{Proceedings of AIIDE}, 2023.

\bibitem[Ji et~al.(2023)]{ji2023survey}
Z.~Ji, N.~Lee, R.~Frieske, T.~Yu, D.~Su, Y.~Xu, E.~Ishii, Y.~J. Bang, A.~Madotto, and P.~Fung.
\newblock Survey of hallucination in natural language generation.
\newblock \emph{ACM Computing Surveys}, 55(12), 2023.

\end{thebibliography}

\end{document}
